% ============================================================================
%  ANEXO 1.A — PLAN INSTITUCIONAL DE CERTIFICACIÓN ISO/IEC 27001:2022
%  audIT Soluciones Tecnológicas SpA - Licitación TFEP-01/2026
% ============================================================================

\section*{Anexo 1.A: Plan institucional de certificación ISO/IEC 27001:2022}\label{sec:1-anexo-iso}
\addcontentsline{toc}{section}{Anexo 1.A: Plan institucional de certificación ISO/IEC 27001:2022}

Conforme a la exigencia de acreditación de Requisito Habilitante del \art{34.1}{22} de las Bases Administrativas TFEP-01/2026, audIT formaliza el siguiente plan institucional vinculante.

\subsection*{1. Individualización del oferente y representación legal}
\begin{itemize}
  \item \textbf{Razón Social:} audIT Soluciones Tecnológicas SpA
  \item \textbf{RUT:} 76.924.310-K
  \item \textbf{Representante Legal:} Alejandro Hermosilla Díaz
  \item \textbf{Cédula de Identidad:} 14.892.341-2
  \item \textbf{Organismo Certificador:} Bureau Veritas Certification S.A. (acreditado ante el Instituto Nacional de Normalización [INN] bajo norma NCh-ISO/IEC 17021, signatario IAF MLA).
  \item \textbf{Norma Internacional:} ISO/IEC 27001:2022 (\textit{Information Security, Cybersecurity and Privacy Protection - Information Security Management Systems}).
  \item \textbf{Expediente de Auditoría Externa:} N.° BV-EXP-2026-CL-8921.
\end{itemize}

\subsection*{2. Alcance del Sistema de Gestión de Seguridad de la Información (SGSI)}
El Sistema de Gestión de Seguridad de la Información de audIT cubre:
\begin{quote}
«El diseño, desarrollo, pruebas, integración, aseguramiento de calidad, implantación y operación de plataformas de software de misión crítica, soluciones de telemetría e Internet de las Cosas (IoT) de borde, y administración de arquitecturas cloud para transporte terrestre, distribución y logística corporativa, conforme a la Declaración de Aplicabilidad (SoA) versión 3.0».
\end{quote}

\subsection*{3. Estado documental del proceso y fundamentación del Plan (\art{34.1}{22})}
audIT acoge formalmente su acreditación a la disposición expresa del \art{34.1}{22} de las Bases Administrativas, que faculta la presentación de un plan institucional de certificación con hitos verificables dentro de los primeros doce meses del Contrato. A la fecha de presentación de esta propuesta técnica, la compañía acredita el siguiente estado documental y procedimental ante Bureau Veritas Certification S.A. (Expediente de Auditoría N.° BV-EXP-2026-CL-8921):
\begin{itemize}
  \item \textbf{Auditoría de Fase 1 (Revisión Documental y Diseño del SGSI):} Concluida y aprobada conforme en noviembre de 2025.
  \item \textbf{Auditoría de Fase 2 (Evaluación de Implementación de 93 Controles):} Ejecutada entre el 10 y el 14 de agosto de 2026 por el equipo auditor senior de Bureau Veritas Certification S.A., con informe técnico favorable y recomendación de certificación con \textbf{cero No Conformidades Mayores}.
  \item \textbf{Estado del Expediente y Tramitación de Emisión:} El expediente BV-EXP-2026-CL-8921 se encuentra en fase final de protocolización ante el Comité de Certificación del organismo acreditador. De conformidad con las Bases, la entrega material del diploma oficial de certificación con acreditación INN/IAF constituye el primer hito vinculante y fiscalizable (\textbf{Hito H1, Mes 1 de Contrato}), respaldado por la declaración jurada notarial del Representante Legal depositada en el Sobre N.º 1.
\end{itemize}

\subsection*{4. Cronograma vinculante de despliegue, protocolización y vigilancia anual}
Para dar cabal cumplimiento a lo dispuesto en el \art{34.1}{22} de las Bases Administrativas (\textit{«plan de certificación con hitos verificables dentro de los primeros doce meses del Contrato»}), en la \tabref{1-iso27001-hitos} se formaliza el cronograma de cumplimiento irrestricto de los cuatro hitos comprometidos.

\begin{tabla}[ancho=completo]{Cronograma vinculante de despliegue, protocolización y vigilancia anual ISO/IEC 27001:2022}{1-iso27001-hitos}{C{15mm} L{32mm} X L{35mm} X}
\textbf{Hito} & \textbf{Plazo Máximo} & \textbf{Actividad y Entregable Verificable} & \textbf{Responsable Corporativo} & \textbf{Criterio de Verificación / Evidencia} \\ \midrule
\textbf{H1} & Mes 1 de Contrato & \textbf{Protocolización y Entrega de Certificado Oficial:} Recepción del diploma emitido por Bureau Veritas e inscripción en registro INN/IAF, entregando copia autorizada a Transportes Curimón S.A. & Representante Legal / Oficial de Seguridad (CISO) & Certificado formal oficial con código de verificación QR y firma digital del organismo acreditador. \\
\textbf{H2} & Mes 4 de Contrato & \textbf{Auditoría Interna de Extensión al Proyecto Curimón:} Evaluación formal de los 93 controles ISO 27001 aplicados a los componentes AKS, pasarelas EdgeHub y base de datos del Caso 10. & Líder de Calidad y Seguridad de audIT SpA & Informe pericial de auditoría interna de proyecto con matriz de riesgos mitigada. \\
\textbf{H3} & Mes 8 de Contrato & \textbf{Revisión por la Dirección y Actualización de SoA:} Evaluación de métricas operacionales de seguridad, análisis de vulnerabilidades OWASP ASVS y actualización de Declaración de Aplicabilidad. & Comité de Seguridad de la Información (CSIC) & Acta formal firmada por Gerencia General y CISO, con matriz de tratamiento de riesgos actualizada. \\
\textbf{H4} & Mes 11 de Contrato & \textbf{Auditoría Externa Anual de Mantenimiento:} Ejecución de la primera auditoría periódica de vigilancia por parte de Bureau Veritas Certification S.A. para certificar la mejora continua. & Bureau Veritas Certification S.A. / CISO & Informe de Auditoría Externa de Vigilancia Anual favorable y sin observaciones bloqueantes. \\
\end{tabla}

Como se desprende de la \tabref{1-iso27001-hitos}, el plan establece controles auditables y verificables en terreno a lo largo de todo el primer año contractual.

\subsection*{5. Declaración formal y firma del representante legal}
En mi calidad de Representante Legal de audIT, declaro bajo fe de juramento que la información expuesta precedentemente es fidedigna y que la empresa asume el compromiso irrestricto de dar cumplimiento cabal a los hitos del presente Plan dentro de los plazos señalados, facultando a Transportes Curimón S.A. para auditar su avance en cualquier momento durante la vigencia del Contrato.

Santiago de Chile, 03 de octubre de 2026.

\vspace{0.8cm}
\begin{center}
  \includegraphics[height=1.5cm]{portadas/activos/media-firma.png}\\[4pt]
  \textbf{Alejandro Hermosilla Díaz}\\
  Representante Legal y Director Ejecutivo\\
  audIT Soluciones Tecnológicas SpA\\
  RUT: 14.892.341-2
\end{center}
